\chapter{The canvas and packed cells}
\label{ch:canvas}

After layout decides where things go, the renderer needs a place to
paint them. That place is the \textbf{canvas}: a 2-D grid of cells
representing the terminal screen. \maya{}'s canvas is unusual in two
ways: each cell is a 64-bit packed integer (not a struct), and the
canvas is double-buffered (front and back) to support the SIMD diff
in the next chapter.

\section{What a cell is}

A terminal cell holds:

\begin{itemize}[leftmargin=*]
  \item A Unicode code point (up to 21 bits).
  \item A foreground colour, a background colour, and a set of text
    attributes.
  \item Optionally, a hyperlink (an OSC-8 URL).
  \item A width marker: 1 for normal characters, 2 for wide CJK and
    emoji, 0 for combining marks and the trailing half of a wide
    cell.
\end{itemize}

A naive struct would look like:

\begin{lstlisting}
struct Cell {
    char32_t cp;        // 4 bytes
    Color    fg;        // 4 bytes
    Color    bg;        // 4 bytes
    uint8_t  attrs;     // 1 byte
    uint8_t  width;     // 1 byte
    uint32_t link;      // 4 bytes
};  // ~18 bytes, padded to 24
\end{lstlisting}

Two problems. First, that is a lot of memory for an 80x24 grid: 1920
cells $\times$ 24 bytes = ~46 KiB, and we need two grids
(double-buffered) = 92 KiB. Second, and more importantly, comparing
two cells means comparing several fields, which is several
instructions and several branches.

\section{Packing}

\maya{} packs every cell into a single 64-bit integer:

\begin{lstlisting}
//  63                                                              0
// +---------+-------+----------+-----+----------------------------+
// | width:2 | tag:2 | hyper:14 | sty:16 |     codepoint:30        |
// +---------+-------+----------+--------+-------------------------+
\end{lstlisting}

The exact bit layout is in \file{maya/include/maya/render/canvas.hpp};
the field names and widths above are illustrative. What matters is
the technique:

\begin{itemize}[leftmargin=*]
  \item \textbf{Code point} occupies the low bits (Unicode is only
    21 bits, so 30 is generous).
  \item \textbf{Style ID} is a 16-bit index into a \code{StylePool}.
    The pool interns full \code{Style} structs (fg, bg, attrs); the
    canvas just stores the ID.
  \item \textbf{Hyperlink ID} likewise indexes a hyperlink pool.
  \item \textbf{Width marker} and a small tag fill the top.
\end{itemize}

Two consequences. First, the canvas shrinks to one quarter of the
naive representation. Second, and crucially: \textbf{cell equality is
\code{uint64\_t} equality}. One machine instruction. No branches. This
is the lever \maya{} uses to diff thousands of cells per microsecond
in the next chapter.

\begin{idea}
Whenever you find yourself comparing two structs in a hot loop, ask
whether you can pack them into a single integer. SIMD becomes
possible the moment your unit of comparison fits in a register lane.
\end{idea}

\section{The style pool}

The 16-bit style ID is what makes packing work. There are typically
only a few dozen distinct styles on screen at any time (text colour,
background, bold, dim, etc.) --- nowhere near the 65536 the ID
supports. The pool is a simple table of \code{Style} structs plus a
hash map from \code{Style} to ID:

\begin{lstlisting}
class StylePool {
public:
    StyleId intern(Style s) {
        auto [it, inserted] = lookup_.try_emplace(s, table_.size());
        if (inserted) {
            table_.push_back(s);
            cache_sgr_.push_back(build_sgr_string(s));
        }
        return it->second;
    }

    const Style&        style(StyleId id) const { return table_[id]; }
    const std::string&  sgr  (StyleId id) const { return cache_sgr_[id]; }

private:
    std::vector<Style>      table_;
    std::vector<std::string> cache_sgr_;     // pre-built SGR per style
    std::unordered_map<Style, StyleId> lookup_;
};
\end{lstlisting}

Two clever bits:

\begin{itemize}[leftmargin=*]
  \item \textbf{Pre-cached SGR.} The first time a \code{Style} is
    interned, the pool builds the ANSI byte string for it
    (\code{ESC[1;38;2;100;180;255m} or whatever the style demands).
    During the render walk, switching style is one \code{memcpy} of
    the cached SGR --- no formatting, no field comparisons.
  \item \textbf{Stable IDs across frames.} The pool persists across
    frames. Frame 2's cells use the same IDs as Frame 1's for the
    same styles, so cell-level equality (which depends on the ID)
    still works.
\end{itemize}

\section{The canvas data structure}

A canvas is two of these grids:

\begin{lstlisting}
class Canvas {
    int width_  = 0;
    int height_ = 0;

    AlignedBuffer front_;   // what's currently on the terminal
    AlignedBuffer back_;    // what we're drawing next

    StylePool    styles_;
    HyperlinkPool links_;

    Rect         damage_;   // dirty region for the diff
    std::vector<Rect> clip_stack_;  // overflow:hidden support
public:
    void clear();
    void set(int x, int y, char32_t cp, StyleId sty);
    void draw_text(int x, int y, std::string_view utf8, StyleId sty);
    void flip();   // swap front and back after diff
    // ...
};
\end{lstlisting}

\code{front\_} is what the user is looking at. \code{back\_} is what
the next frame should look like. Each frame:

\begin{enumerate}[leftmargin=*]
  \item Clear \code{back\_} (zero-fill, AVX-512 \code{vmovdqu64} when
    available).
  \item Walk the element tree; for each visible cell, call
    \code{back\_.set(x, y, cp, sty)}.
  \item Diff \code{front\_} against \code{back\_}; emit ANSI for the
    changed cells.
  \item Swap roles --- \code{back\_} becomes \code{front\_}.
\end{enumerate}

This is the classic \emph{double-buffered} pattern.

\subsection{Why 64-byte alignment?}

The \code{AlignedBuffer} class allocates cells aligned to 64-byte
boundaries (a full CPU cache line). The diff loops in the next chapter
load 256 bits at a time (AVX2) or 512 bits (AVX-512); unaligned loads
cost cycles when they cross a cache line. Aligning to 64 bytes
guarantees the SIMD loads never split a line.

\section{Drawing text}

The text-drawing path looks innocuous but does Unicode work:

\begin{lstlisting}
void Canvas::draw_text(int x, int y, std::string_view utf8, StyleId sty) {
    int cx = x;
    for (char32_t cp : utf8_decode(utf8)) {
        int w = unicode_width(cp);
        if (w == 0) {
            // combining mark - attach to previous cell (skipped for brevity)
        } else if (w == 1) {
            set(cx, y, cp, sty);
            cx += 1;
        } else if (w == 2) {
            set(cx,     y, cp, sty, /*width=*/2);
            set(cx + 1, y, 0,  sty, /*width=*/0);  // trailing half
            cx += 2;
        }
    }
}
\end{lstlisting}

The wide-character handling is what keeps CJK and emoji from
corrupting the layout. The lookup table at
\file{maya/include/maya/text/unicode\_width\_table.hpp} is generated
from the Unicode standard's \code{EastAsianWidth.txt} and
\code{emoji-data.txt} by the script
\file{maya/scripts/gen\_unicode\_width.py}.

\section{The damage rectangle}

To avoid diffing the entire canvas every frame, \maya{} tracks the
bounding box of any cell that has been written this frame:

\begin{lstlisting}
void Canvas::set(int x, int y, char32_t cp, StyleId sty) {
    back_[y * width_ + x] = pack(cp, sty);
    damage_ = damage_ | Rect{x, y, 1, 1};   // union the rectangle
}
\end{lstlisting}

The diff in the next chapter visits only the cells inside
\code{damage\_}. If most of the screen is unchanged (the typical
case for steady-state UIs), the diff is fast in proportion to what
moved, not in proportion to the screen size.

\subsection{The clip stack}

A box with \code{overflow: Hidden} (a scrollable list, a modal) needs
its rendering clipped to its rectangle. \maya{} pushes the clip
rectangle onto a stack and intersects every \code{set()} with the
current top:

\begin{lstlisting}
void Canvas::push_clip(Rect r) { clip_stack_.push_back(intersect(top(), r)); }
void Canvas::pop_clip()        { clip_stack_.pop_back(); }
\end{lstlisting}

Children that try to draw outside the clip simply have their writes
silently dropped.

\section{Resize}

When \code{SIGWINCH} fires, the runtime queries the new size and
calls \code{canvas.resize(w, h)}. The buffers reallocate (preserving
old contents is unnecessary --- the next frame redraws everything),
the damage rect is set to the full new screen, and the next paint
emits ANSI for the whole grid.

\subsection{Shrinking after a wide window}

A neat detail: if the user grows the terminal to 200 cells wide, runs
the app, then shrinks back to 80 cells, the AlignedBuffer would
otherwise stay at the wide capacity forever. \maya{} ships
\code{shrink\_to\_fit\_if\_oversized} that releases capacity when at
most half of it is in use, to avoid permanent over-allocation but
also to avoid thrashing the allocator on tiny oscillations.

\section{Putting it together}

A frame, end to end, in pseudocode:

\begin{lstlisting}
void Renderer::frame(const Element& root) {
    canvas.clear_back();
    layout::compute(layout_tree, /*root=*/0, canvas.width());
    paint_visit(canvas, layout_tree, /*node=*/0);   // fills back_

    std::string ansi;
    diff(canvas.front(), canvas.back(),
         canvas.damage(), canvas.styles(), ansi);   // SIMD-accelerated
    writer.write(ansi);

    canvas.flip();
}
\end{lstlisting}

The chapter has set up everything but the diff. That is next.

\begin{recap}
The canvas is a double-buffered grid of packed 64-bit cells.
Packing makes cell equality a single integer comparison; the style
pool interns full \code{Style} structs into 16-bit IDs. A damage
rectangle limits the next chapter's diff to the changed region.
Wide characters and clip regions are handled in the same data
structure.
\end{recap}

\section*{Workshop}

\begin{enumerate}[leftmargin=*]
  \item Open \file{maya/include/maya/render/canvas.hpp} and find the
    \code{pack()} / \code{unpack()} helpers. Note the bit layout.
  \item Run an example with wide characters
    (\file{examples/messenger.cpp}). Notice how alignment stays
    correct when the text contains emoji.
  \item Estimate, on the back of an envelope, the byte cost of an
    80x24 canvas (1920 cells), then double it for the two buffers,
    then add typical pool sizes. You should land near 32 KiB ---
    well under L1 cache.
\end{enumerate}
